\documentclass[aspectratio=169,11pt]{beamer}
% Shared Beamer style for practice contest solution walkthroughs.
\usepackage[utf8]{inputenc}
\usepackage[T1]{fontenc}
\usepackage{lmodern}
\usepackage{amsmath,amssymb}
\usepackage{xcolor}
\usepackage{hyperref}
\usepackage{listings}

\definecolor{CPNavy}{HTML}{172033}
\definecolor{CPBlue}{HTML}{2563EB}
\definecolor{CPGreen}{HTML}{16A34A}
\definecolor{CPOrange}{HTML}{F97316}
\definecolor{CPGray}{HTML}{64748B}
\definecolor{CPLight}{HTML}{F8FAFC}
\definecolor{CPCodeBg}{HTML}{F1F5F9}
\definecolor{CPCodeRule}{HTML}{CBD5E1}

\usetheme{default}
\usefonttheme{professionalfonts}
\setbeamertemplate{navigation symbols}{}
\setbeamersize{text margin left=0.65cm,text margin right=0.65cm}
\setbeamertemplate{itemize item}{\color{CPBlue}$\blacktriangleright$}
\setbeamertemplate{itemize subitem}{\color{CPGray}$\triangleright$}

\setbeamercolor{normal text}{fg=CPNavy,bg=white}
\setbeamercolor{frametitle}{fg=white,bg=CPNavy}
\setbeamercolor{title}{fg=white,bg=CPNavy}
\setbeamercolor{subtitle}{fg=CPLight,bg=CPNavy}
\hypersetup{colorlinks=true,urlcolor=CPBlue}

\setbeamertemplate{frametitle}{%
  \nointerlineskip%
  \begin{beamercolorbox}[wd=\paperwidth,ht=0.88cm,dp=0.22cm,leftskip=0.55cm,rightskip=0.35cm]{frametitle}
    \usebeamerfont{frametitle}\insertframetitle\hfill\small\insertframenumber/\inserttotalframenumber
  \end{beamercolorbox}%
}

\setbeamertemplate{footline}{%
  \leavevmode%
  \hbox{%
    \begin{beamercolorbox}[wd=.58\paperwidth,ht=2.4ex,dp=1ex,leftskip=.5cm]{author in head/foot}%
      \color{CPGray}\insertshorttitle
    \end{beamercolorbox}%
    \begin{beamercolorbox}[wd=.42\paperwidth,ht=2.4ex,dp=1ex,rightskip=.5cm plus1fil]{date in head/foot}%
      \hfill\color{CPGray}\insertshortauthor
    \end{beamercolorbox}%
  }%
}

\setbeamertemplate{title page}{%
  \vspace*{1.25cm}
  \begin{beamercolorbox}[wd=\paperwidth,sep=0.65cm,leftskip=0.15cm,rightskip=0.15cm]{title}
    {\color{white}\Huge\bfseries\inserttitle\par}
    \vspace{0.35cm}
    {\color{CPLight}\Large\insertsubtitle\par}
    \vspace{0.6cm}
    {\color{CPLight}\large\insertauthor\par}
  \end{beamercolorbox}
  \vfill
}

\newcommand{\sectiontag}[1]{\textbf{\color{CPBlue}#1}}
\newenvironment{tightitemize}{\begin{itemize}\small\setlength{\itemsep}{0.18cm}}{\end{itemize}}
\lstdefinestyle{cpcode}{
  language=Python,
  basicstyle=\ttfamily\tiny,
  keywordstyle=\color{CPBlue}\bfseries,
  commentstyle=\color{CPGray}\itshape,
  stringstyle=\color{CPGreen},
  backgroundcolor=\color{CPCodeBg},
  frame=single,
  framerule=0.25pt,
  rulecolor=\color{CPCodeRule},
  showstringspaces=false,
  columns=fullflexible,
  keepspaces=true,
  breaklines=true,
  breakatwhitespace=false,
  tabsize=4,
  xleftmargin=0.08cm,
  xrightmargin=0.08cm,
  aboveskip=0.08cm,
  belowskip=0.08cm
}


\title[fenwick]{Fenwick Tree}
\subtitle{Day 1 solution walkthrough}
\author[Solution Deck]{Competitive Programming Bootcamp}
\date{}

\begin{document}

\begin{frame}[plain]
  \titlepage
\end{frame}

% BEGIN GENERATED STATEMENT INTERPRETATION
\begin{frame}{Interpret the Word Problem}
\small
\sectiontag{Statement excerpts:}
\begin{tightitemize}
  \item \textit{``updates and prefix sum queries''}
  \item \textit{``a query for the value of \$a[0] + a[1] + \textbackslash{}ldots + a[i-1]\$''}
\end{tightitemize}

\vspace{0.18cm}
\sectiontag{Programming interpretation:}
\begin{tightitemize}
  \item The array starts at zero and receives point increments.
  \item Every query asks for a prefix sum, not a closed interval ending at i.
  \item A Fenwick tree gives both operations logarithmic time.
\end{tightitemize}
\end{frame}
% END GENERATED STATEMENT INTERPRETATION







\begin{frame}{Problem Model}
\small
\sectiontag{Textbook connection:} CP4 2.4.3 Fenwick (Binary Indexed) Tree

\vspace{0.25cm}
\sectiontag{Core technique:} Fenwick tree for dynamic prefix sums

\vspace{0.25cm}
\begin{tightitemize}
  \item The array starts with all zeroes.
  \item '+ i v' adds v to one position.
  \item '? i' asks for the sum of positions before i.
\end{tightitemize}
\end{frame}

\begin{frame}{How to Recognize the Approach}
\begin{tightitemize}
  \item Naive prefix queries are too slow after many updates.
  \item Fenwick trees support point updates and prefix sums in logarithmic time.
  \item The main implementation detail is converting 0-based problem indices to 1-based tree indices.
\end{tightitemize}
\end{frame}

\begin{frame}{Algorithm}
\begin{tightitemize}
  \item Build a Fenwick tree array ft[1..N].
  \item update(i, v): add v to ft[i], then jump i += lsone(i).
  \item query(1, j): accumulate ft[j], then jump j -= lsone(j).
  \item For '+ i v', call update(i + 1, v).
  \item For '? i', query positions 1 through i.
\end{tightitemize}
\end{frame}

\begin{frame}{Walkthrough of the Key State}
\begin{tightitemize}
  \item lsone(x) = x \& -x identifies the size of the chunk controlled by x.
  \item Update moves upward to every chunk containing the changed index.
  \item Query moves downward through disjoint chunks that cover the prefix.
\end{tightitemize}
\end{frame}

\begin{frame}{Why the Algorithm Is Correct}
\begin{tightitemize}
  \item Each Fenwick cell stores the sum of a fixed suffix chunk ending at that index.
  \item The update visits exactly the chunks that include the changed position.
  \item The query decomposes a prefix into disjoint stored chunks, so their sum equals the prefix sum.
\end{tightitemize}
\end{frame}

\begin{frame}{Complexity and Implementation Notes}
\small
\sectiontag{Complexity:} O(log N) per update/query and O(N) memory.

\vspace{0.3cm}
\begin{tightitemize}
  \item Use integer arithmetic; sums may exceed small integer ranges in other languages.
  \item Kattis '? i' is exclusive of i, which is why query(1, i) is correct after 1-based conversion.
\end{tightitemize}
\end{frame}



% BEGIN GENERATED CODE WALKTHROUGH
\begin{frame}[fragile]{Code: Build Tree Storage}
\scriptsize
\sectiontag{Source lines:} \texttt{fenwick.py:34--47}

\vspace{0.08cm}
\begin{columns}[T,totalwidth=\textwidth]
\begin{column}{0.58\textwidth}
\begin{lstlisting}[style=cpcode]
class FTree:
    def __init__(self, f):
        self.n = len(f)
        self.ft = [0] * (self.n + 1)

        for i in range(1, self.n + 1):
            self.ft[i] += f[i - 1]

            parent = i + self.lsone(i)

            if parent <= self.n:
                self.ft[parent] += self.ft[i]
\end{lstlisting}
\end{column}
\begin{column}{0.38\textwidth}
\sectiontag{What this chunk does}
\begin{tightitemize}
  \item The Fenwick array has one extra cell because the structure is 1-based.
  \item The build loop pushes each value into the parent responsible for a larger range.
  \item For this problem the initial array is all zeros, but the constructor is general.
\end{tightitemize}
\end{column}
\end{columns}
\end{frame}

\begin{frame}[fragile]{Code: Prefix and Range Sum}
\scriptsize
\sectiontag{Source lines:} \texttt{fenwick.py:48--65}

\vspace{0.08cm}
\begin{columns}[T,totalwidth=\textwidth]
\begin{column}{0.58\textwidth}
\begin{lstlisting}[style=cpcode]
def lsone(self, s):
    return s & (-s)

def query(self, i, j):
    if i > 1:
        return self.query(1, j) - self.query(1, i - 1)

    s = 0

    while j > 0:
        s += self.ft[j]
        j -= self.lsone(j)

    return s
\end{lstlisting}
\end{column}
\begin{column}{0.38\textwidth}
\sectiontag{What this chunk does}
\begin{tightitemize}
  \item lsone isolates the size of the range represented by the current tree node.
  \item A prefix query repeatedly jumps to the previous responsible node.
  \item A range query is handled as the difference of two prefix sums.
\end{tightitemize}
\end{column}
\end{columns}
\end{frame}

\begin{frame}[fragile]{Code: Point Updates}
\scriptsize
\sectiontag{Source lines:} \texttt{fenwick.py:67--90}

\vspace{0.08cm}
\begin{columns}[T,totalwidth=\textwidth]
\begin{column}{0.58\textwidth}
\begin{lstlisting}[style=cpcode]
def update(self, i, v):
    while i <= self.n:
        self.ft[i] += v
        i += self.lsone(i)

def select(self, k):
    p = 1

    while (p * 2) <= self.n:
        p *= 2

    i = 0

    while p > 0:
        if i + p <= self.n and k > self.ft[i + p]:
            k -= self.ft[i + p]
            i += p

        p //= 2

    return i + 1
\end{lstlisting}
\end{column}
\begin{column}{0.38\textwidth}
\sectiontag{What this chunk does}
\begin{tightitemize}
  \item Updating position i means adding the delta to every Fenwick node whose range contains i.
  \item The index moves upward by lsone(i).
  \item select is library functionality and is not needed by this Kattis problem.
\end{tightitemize}
\end{column}
\end{columns}
\end{frame}

\begin{frame}[fragile]{Code: Apply Updates}
\scriptsize
\sectiontag{Source lines:} \texttt{fenwick.py:93--112}

\vspace{0.08cm}
\begin{columns}[T,totalwidth=\textwidth]
\begin{column}{0.58\textwidth}
\begin{lstlisting}[style=cpcode]
def main():
    input = sys.stdin.buffer.readline

    n, q = map(int, input().split())

    ft = FTree([0] * n)

    output = []

    for _ in range(q):
        command = input().split()

        if command[0] == b"+":
            index = int(command[1])
            value = int(command[2])

            ft.update(index + 1, value)
\end{lstlisting}
\end{column}
\begin{column}{0.38\textwidth}
\sectiontag{What this chunk does}
\begin{tightitemize}
  \item The Kattis indices are 0-based, so updates use index + 1 inside the tree.
  \item A plus command increments one array position.
  \item The Fenwick update pushes the delta to every responsible tree node.
\end{tightitemize}
\end{column}
\end{columns}
\end{frame}

\begin{frame}[fragile]{Code: Answer Prefix Queries}
\scriptsize
\sectiontag{Source lines:} \texttt{fenwick.py:114--125}

\vspace{0.08cm}
\begin{columns}[T,totalwidth=\textwidth]
\begin{column}{0.58\textwidth}
\begin{lstlisting}[style=cpcode]
        else:
            index = int(command[1])

            output.append(str(ft.query(1, index)))

    sys.stdout.write("\n".join(output))

main()
\end{lstlisting}
\end{column}
\begin{column}{0.38\textwidth}
\sectiontag{What this chunk does}
\begin{tightitemize}
  \item The query '? i' maps directly to tree positions 1 through i.
  \item Each prefix sum is appended to the output buffer.
  \item All answers are printed once at the end.
\end{tightitemize}
\end{column}
\end{columns}
\end{frame}
% END GENERATED CODE WALKTHROUGH

\begin{frame}{Source}
\small
\sectiontag{Solution file:} \texttt{practice\_contests/day\_1/fenwick.py}

\vspace{0.3cm}
\sectiontag{Problem:} \url{https://open.kattis.com/problems/fenwick}

\vspace{0.45cm}
Use this deck with the source code open beside it. The slides explain the reasoning; the code shows the exact input handling and output formatting.
\end{frame}

\end{document}
